% Slovak version
\section*{Anotácia}

\spacenewline
\begin{tabular}{@{}p{0.23\linewidth}p{0.72\linewidth}}
\multicolumn{2}{@{}l}{Slovenská technická univerzita v Bratislave} \\
\multicolumn{2}{@{}l}{FAKULTA INFORMATIKY A INFORMAČNÝCH TECHNOLÓGIÍ} \\
Študijný program: & Aplikovaná Informatika \\
Autor: & Ing. Viktor Valaštín \\
Dizertačná práca: & Internet hodnôt \\
Vedúci práce: & Prof. Ivan Kotuliak, PhD \\
Konzultant: & Ing. Kristián Košťál, PhD \\
May 2026
\end{tabular}

\spacedoublenewline

\spacenewline

Internet vecí ako veľmi známy pojem pokrýva schopnosť zariadení zdieľať informácie medzi sebou a interagovať cez internet.
Internet hodnôt je evolúciou Internetu vecí, ktorá entitám umožňuje zdieľať hodnotu.
S vývojom internetového bankovníctva a gig ekonomiky sa hodnota zdieľaná cez internet enormne zvýšila.
Aby sa dosiahol skutočný potenciál Internetu hodnôt, potrebujeme dôveryhodný a bezpečný spôsob výmeny hodnôt medzi entitami.
Využitie technológie decentralizovaných záznamov (decentralized ledger technology) môže poskytnúť stabilný a bezpečný spôsob vytvárania a výmeny hodnôt.
Dizertačná práca sa zameriava na definovanie Internetu hodnôt v kontexte blockchainu a implementáciu sady nástrojov, ktoré umožnia nové spôsoby výmeny hodnôt.
Ďalej práca pokrýva zdieľanie hodnôt s autonómnymi agentami založenými na umelej inteligencii ako perspektívny nástroj na prácu s aktívami.
Kľúčovou inováciou práce je návrh viacvrstvovej agentickej architektúry a formálneho rámca pre viacrozmernú valuáciu aktív.
Predstavené nástroje sú navrhnuté tak, aby boli využiteľné s virtuálnymi aj fyzickými aktívami.



% English version
\section*{Annotation}

\spacenewline
\begin{tabular}{@{}p{0.23\linewidth}p{0.72\linewidth}}
\multicolumn{2}{@{}l}{Slovak University of Technology Bratislava} \\
\multicolumn{2}{@{}l}{FACULTY OF INFORMATICS AND INFORMATION TECHNOLOGIES} \\
Study program: & Applied Informatics \\
Author: & Ing. Viktor Valaštín \\
Dissertation thesis: & Internet of Value \\
Supervisor: & Prof. Ivan Kotuliak \\
Consultant: & Dr. Kristián Košťál \\
May 2026
\end{tabular}

\spacedoublenewline

Internet of Things, as a well-known term, covers the ability of devices
to share information and interact with each other over the Internet.
The Internet of Value is an evolution of IoT in which entities can share value.
With the advent of internet banking and the gig economy, the value shared over the Internet has increased significantly.
To realise the potential of the Internet of Value, entities need a trustless and secure way to
exchange value. Decentralized ledger technology can provide a stable and secure foundation for creating and exchanging value. This dissertation defines the Internet of Value in the context of blockchain and implements a set of tools
that enable new forms of value exchange. Furthermore, the thesis covers value sharing with autonomous agents based on artificial intelligence as a prospective tool for asset management. The key novelty of the thesis is the design of a multi-layered agentic architecture and a formal framework for multidimensional asset valuation. The presented tools
are designed for use with both virtual and physical assets.
